\section{Limitations}
\label{sec:limitations}

The paper makes a bounded systems claim: backend choice should depend on checkpoint format and model scale. It does not prove a universal selector, universal thresholds, or universal superiority of any I/O API.

The primary measurements come from Modal H100 containers under gVisor, where \texttt{io\_uring} is unavailable. This restricts the evaluation to sync and async backends; the \texttt{io\_uring} backend, which showed significant advantages on bare-metal H100 with local NVMe in prior work, cannot be evaluated in this environment. Practitioners with bare-metal access and working \texttt{io\_uring} may observe a three-way regime structure rather than the two-way comparison reported here.

The model ladder covers six public checkpoints from 0.6B to 32B parameters and two layouts. It does not cover GGUF, quantized-only deployments, object stores, distributed filesystems, RAID variants, direct-I/O configurations, Windows/macOS backends, or multi-node distributed loading. ``Load time'' means different things at different layers: Rust timings approximate backend-attributed disk-to-tensor work; vLLM timings add Python integration overhead, engine initialization, and CUDA setup. A Rust result should not be quoted as a serving result.

The results compare two specific implementations (sync POSIX and Tokio async), not the best possible implementation of each API. A differently tuned sync backend with adaptive parallelism or a differently configured async runtime might shift the exact crossover points. This does not invalidate the central claim that checkpoint format changes the optimal backend, because that claim depends on the workload-shape explanation (contiguous vs.\ partitioned reads), not on any single implementation being universally optimal.

Modal containers introduce CDN-induced variance from model downloads, which inflates measurement ranges relative to local-NVMe environments. The five-repetition statistical methodology accounts for this by reporting medians and ranges rather than single-point estimates, but individual-cell precision is lower than on local storage.

The Python integration layer is not independently benchmarked. The Rust-to-vLLM stack provides end-to-end evidence that backend effects propagate through the serving stack, but the Python-level contribution is not separately quantified.
